\par\Needspace{6\baselineskip}\medskip
\begin{center}\hypertarget{native-169-article122}{}\pdfbookmark[2]{122.}{article-122}122.\end{center}
Quum incognitarum valores in hypothesi secunda et tertia supponendi
quodammodo arbitrarii sint, si modo ab hypothesi prima non nimis differant,
praetereaque caveatur, ne ratio $(a''-a):(b''-b)$ ad aequalitatem huius $(a'-a):(b'-b)$ convergat, plerumque statui solet $a'=a$, $b''=b$. Duplex hinc lucrum
derivatur: namque non solum formulae pro $\xi,\eta$ paullo adhuc simpliciores evadunt, sed pars quoque calculi primi eadem manebit in hypothesi secunda,
aliaque pars in tertia.

Est tamen casus, ubi aliae rationes ab hac consuetudine discedere suadent: fingamus enim, $X$ habere formam $X'-x$, atque $Y$ hanc $Y'-y$, functionesque $X',Y'$ per problematis naturam ita comparatas esse, ut erroribus
mediocribus in valoribus ipsarum $x,y$ commissis perparum afficiantur, sive ut
$\left(\dfrac{\mathrm{d}X'}{\mathrm{d}x}\right)\mathpunct{\raisebox{.45ex}{,}}$ $\left(\dfrac{\mathrm{d}X'}{\mathrm{d}y}\right)\mathpunct{\raisebox{.45ex}{,}}$ $\left(\dfrac{\mathrm{d}Y'}{\mathrm{d}x}\right)\mathpunct{\raisebox{.45ex}{,}}$ $\left(\dfrac{\mathrm{d}Y'}{\mathrm{d}y}\right)$ sint quantitates perexiguae, patetque, differentias
\SourceLeaf{170}inter valores istarum functionum systemati $x=\xi$, $y=\eta$ respondentes, eosque
qui ex $x=a$, $y=b$ prodeunt, ad ordinem quasi altiorem referri posse, quam
differentias $\xi-a$, $\eta-b$; at valores illi sunt $X'=\xi$, $Y'=\eta$, hi vero $X'=a+A$, $Y'=b+B$, unde sequitur, $a+A$, $b+B$ esse valores multo exactiores
ipsarum $x,y$, quam $a,b$. Quibus si hypothesis secunda superstruitur, persaepe
aequationibus $X=0$, $Y=0$ tam exacte iam satisfit, ut ulterius progredi haud
opus sit; sin secus, eodem modo ex hypothesi secunda tertia formabitur faciendo
$a''=a'+A'=a+A+A'$, $b''=b'+B'=b+B+B'$, unde tandem, si
nondum satis praecisa reperitur, quarta ad normam art. 120 elicietur.
